\chapter{Rectal Cancer Surgery}

\section*{What Is This Surgery?}
Rectal cancer surgery, commonly referred to as radical proctectomy, encompasses a range of surgical interventions aimed at excising cancerous lesions located in the rectum, which is the terminal segment of the large intestine. This surgery is primarily indicated for rectal adenocarcinoma, which typically arises within the distal 15 cm of the rectum. The procedure is critical for achieving curative outcomes in patients with resectable tumors, with the overarching goal of complete (R0) resection while preserving anal sphincter function whenever feasible. Depending on the tumor's location and stage, surgical options may include low anterior resection (LAR) with anastomosis or abdominoperineal resection (APR) resulting in a permanent colostomy. Advances in surgical techniques, particularly minimally invasive approaches such as laparoscopic and robotic-assisted surgeries, have significantly improved patient recovery times and reduced postoperative complications.

\section*{Alternative Names}
\begin{itemize}
  \item Proctectomy
  \item Total Mesorectal Excision (TME)
  \item Low Anterior Resection (LAR)
  \item Ultra-Low Anterior Resection (ULAR)
  \item Abdominoperineal Resection (APR)
  \item Intersphincteric Resection (ISR)
\end{itemize}

\section*{Relevant Surgical Anatomy}
The rectum is a muscular tube approximately 12-15 cm long, extending from the rectosigmoid junction at the level of the S3 vertebra to the anal canal. It is characterized by the absence of teniae coli, haustra, and epiploic appendages, which are present in the colon. The rectum is supported inferiorly by the levator ani muscles and is closely related to various pelvic organs: anteriorly to the bladder and prostate in males, and to the vagina and uterus in females. 

A key anatomical feature in rectal cancer surgery is the mesorectum, a fatty tissue envelope that surrounds the rectum and contains vital structures such as blood vessels, lymphatics, and nerves. The mesorectal fascia, which delineates the mesorectum, is crucial for surgical dissection, allowing for the removal of the rectum and surrounding tissues while minimizing blood loss and preserving nerve function.

The arterial supply to the rectum is primarily provided by the superior rectal artery (a continuation of the inferior mesenteric artery), with additional contributions from the middle rectal arteries (branches of the internal iliac arteries) and inferior rectal arteries (branches of the internal pudendal arteries). Venous drainage parallels the arterial supply, with the superior rectal vein draining into the inferior mesenteric vein and the middle and inferior rectal veins draining into the internal iliac veins. Lymphatic drainage follows the superior rectal vessels to the inferior mesenteric lymph nodes and laterally to the internal iliac nodes, particularly for lower rectal tumors. The pelvic autonomic nerves, including the hypogastric and pelvic splanchnic nerves, form the inferior hypogastric plexus, which is critical for maintaining bladder, bowel, and sexual function. Surgical landmarks include the sacral promontory, the ureters crossing the iliac vessels, and the levator ani muscles.

\section*{Surgical Principle \& Rationale}
The cornerstone of rectal cancer surgery is the principle of Total Mesorectal Excision (TME), a technique pioneered by Professor R.J. Heald in the 1980s. TME involves the meticulous dissection of the rectum along with the surrounding mesorectum as a single, intact specimen. This approach aims to achieve complete excision of the primary tumor and its lymphatic drainage, thereby maximizing the chances of obtaining clear surgical margins (R0 resection) and minimizing the risk of local recurrence.

The rationale behind TME is rooted in the embryological development of the rectum, which is encased in a distinct fascial layer. Dissecting along this avascular plane reduces intraoperative blood loss, preserves critical autonomic nerves, and ensures the integrity of the mesorectal package. The specific type of TME performed—whether partial, total, or ultralow—depends on the tumor's distance from the anal verge. For higher rectal tumors, a partial mesorectal excision may suffice, while lower tumors necessitate a total or ultralow TME to ensure adequate distal margins and comprehensive lymph node removal. The ultimate technical goal is to balance oncological radicality with functional preservation, striving for sphincter-sparing procedures whenever possible without compromising cancer cure.

\section*{History \& Surgical Evolution}
The history of rectal cancer surgery has evolved significantly over the past century. Early surgical interventions in the 19th century were often associated with high morbidity and mortality, with many procedures being palliative rather than curative. The modern era began in 1908 when Ernest Miles introduced the abdominoperineal resection (APR), a radical approach that involved the removal of the rectum and anus through both abdominal and perineal incisions, resulting in a permanent colostomy. While effective for local control, this technique was associated with substantial morbidity and a profound impact on patients' quality of life.

In the mid-20th century, the focus shifted towards sphincter-preserving techniques. Claude F. Dixon's introduction of the anterior resection in 1939 allowed for anastomosis and avoided a permanent stoma for higher rectal tumors. However, local recurrence rates remained a significant challenge, often exceeding 30-50\%. The landscape transformed in 1982 with R.J. Heald's introduction of TME, which emphasized the complete removal of the mesorectum as a single unit. This technique dramatically reduced local recurrence rates to less than 10\%, as evidenced by the landmark Dutch TME trial. The late 20th and early 21st centuries saw the advent of laparoscopic and robotic approaches, further enhancing patient recovery and minimizing surgical trauma while maintaining excellent oncological outcomes.

\section*{Clinical Indications}
Rectal cancer surgery is indicated in several clinical scenarios, including:

\begin{itemize}
  \item Curative resection of primary, resectable rectal adenocarcinoma (Stage I, II, or III).
  \item Locally advanced rectal cancers (Stage II or III) following neoadjuvant chemoradiotherapy to downstage the tumor and improve R0 resection rates.
  \item Symptomatic relief of obstructing rectal tumors, even in palliative settings, to enhance quality of life.
  \item Management of significant rectal bleeding or perforation caused by a rectal tumor unresponsive to conservative measures.
  \item Salvage surgery for select cases of recurrent rectal cancer where local resection is feasible and offers a chance for cure or significant palliation.
  \item Surgical intervention for benign conditions such as refractory inflammatory bowel disease with severe dysplasia or familial adenomatous polyposis at high risk of malignancy.
  \item Surgical excision of neuroendocrine tumors of the rectum that are amenable to resection.
\end{itemize}

\section*{Clinical Positioning}
Rectal cancer surgery is primarily a first-line curative treatment for patients diagnosed with resectable rectal adenocarcinoma. For early-stage tumors (Stage I, T1-T2 N0 M0), surgery alone may suffice. However, for locally advanced tumors (Stage II and III, T3-T4 or N+), surgery is often performed as the definitive treatment following neoadjuvant chemoradiotherapy. This preoperative therapy aims to shrink the tumor, sterilize microscopic disease, and improve the likelihood of a complete surgical resection with negative margins, thereby enhancing the curative potential of the surgery. While primarily a first-line approach, surgery can also serve as a secondary or salvage procedure in cases of local recurrence after initial non-operative management or for palliative purposes to manage complications such as obstruction, bleeding, or perforation in patients with unresectable metastatic disease, where the goal is symptom relief rather than cure.

\section*{Surgical Technique \& Procedure}
Rectal cancer surgery is performed under general anesthesia. The patient is typically positioned in a modified lithotomy position to facilitate both abdominal and perineal access, particularly for abdominoperineal resections or ultralow anterior resections. The surgical approach may be open (via a midline laparotomy), laparoscopic, or robotic, with minimally invasive techniques increasingly preferred due to their associated benefits of reduced recovery times and postoperative pain.

The key operative steps generally include:

\begin{itemize}
  \item \textbf{Abdominal Exploration and Mobilization:} Initial assessment for metastatic disease in the liver, peritoneum, and regional lymph nodes. Mobilization of the left colon and splenic flexure is performed to ensure adequate length for a tension-free anastomosis, which is crucial for preventing anastomotic leaks.
  
  \item \textbf{Vascular Ligation:} The inferior mesenteric artery (IMA) and vein (IMV) are ligated at their origin (high ligation) or more distally (low ligation), depending on oncological principles and the need for collateral blood supply to the remaining colon. This ensures adequate oncological clearance of lymph nodes along the IMA and controls the primary blood supply to the resected segment.
  
  \item \textbf{Total Mesorectal Excision (TME):} This critical step involves meticulous sharp dissection in the avascular plane between the mesorectum and the pelvic sidewall fascia (presacral fascia posteriorly, Denonvilliers' fascia anteriorly in males, rectovaginal septum in females). The dissection proceeds circumferentially down to the pelvic floor, ensuring the entire mesorectum is removed intact as a single specimen, thereby minimizing tumor cell spillage and achieving clear circumferential margins.
  
  \item \textbf{Rectal Transection:} The rectum is transected proximally at a healthy margin (typically 5 cm from the tumor for higher lesions) and distally, with the distal margin dictated by the tumor's location and the feasibility of sphincter preservation. For Low Anterior Resection (LAR), a linear stapler is often used for distal transection, aiming for a 1-2 cm distal margin from the tumor for TME.
  
  \item \textbf{Anastomosis or Colostomy:}
    \begin{itemize}
      \item For Low Anterior Resection (LAR), the remaining colon is brought down into the pelvis and anastomosed to the rectal stump or anal canal, typically using a circular stapler to create a colorectal or coloanal anastomosis.
      \item For Abdominoperineal Resection (APR), the anus and distal rectum are removed through a separate perineal incision, and a permanent end colostomy is created in the left lower abdominal wall.
    \end{itemize}
  
  \item \textbf{Diverting Stoma (Optional but common):} For very low anastomoses (especially coloanal) or in patients with significant risk factors for anastomotic leak, a temporary diverting loop ileostomy or colostomy is often created to protect the anastomosis and reduce the consequences of a potential leak. This stoma is typically reversed 8-12 weeks later.
  
  \item \textbf{Closure:} The abdominal incision is closed in layers, and any perineal wound is closed. Pelvic drains may be placed to monitor for bleeding or anastomotic leakage, though their routine use is debated.
\end{itemize}

\section*{Preoperative Workup \& Preparation}
Thorough preoperative preparation is crucial for optimizing outcomes in rectal cancer surgery, minimizing complications, and ensuring patient readiness for a major operation. This typically involves a comprehensive medical evaluation to assess the patient's overall health and fitness for major surgery.

Key preparatory steps include:

\begin{itemize}
  \item \textbf{Diagnostic Workup:} Detailed history and physical examination, complete blood count (CBC), electrolyte panel, liver function tests (LFTs), renal function tests, and carcinoembryonic antigen (CEA) levels.
  
  \item \textbf{Staging Imaging:} High-resolution pelvic MRI is essential for local staging (tumor depth, nodal involvement, relationship to mesorectal fascia, extramural vascular invasion) and guiding neoadjuvant therapy. A CT scan of the chest, abdomen, and pelvis is performed for distant metastatic staging.
  
  \item \textbf{Neoadjuvant Therapy:} For locally advanced tumors (Stage II/III), completion of preoperative chemoradiotherapy, followed by a restaging evaluation, is essential. A typical interval of 8-12 weeks post-chemoradiotherapy is observed before surgery to allow for maximal tumor regression and tissue recovery.
  
  \item \textbf{Anesthesia and Cardiac Clearance:} Comprehensive cardiac and pulmonary evaluation, including ECG, chest X-ray, and potentially echocardiogram or stress testing, to ensure the patient can tolerate general anesthesia and the physiological stress of surgery.
  
  \item \textbf{Bowel Preparation:} Mechanical bowel preparation (e.g., polyethylene glycol solution) combined with oral antibiotics (e.g., neomycin and metronidazole) is often prescribed to reduce the bacterial load in the colon and decrease the risk of surgical site infection and anastomotic complications.
  
  \item \textbf{Medication Review:} Adjustment or discontinuation of certain medications, particularly anticoagulants or antiplatelet agents, in consultation with the surgical team and cardiology.
  
  \item \textbf{NPO Status:} Patients are typically instructed to be NPO (nil per os) for a specified period (usually 6-8 hours for solids, 2 hours for clear liquids) before surgery.
  
  \item \textbf{Prophylactic Antibiotics:} Intravenous broad-spectrum antibiotics are administered shortly before incision to prevent surgical site infections.
  
  \item \textbf{Informed Consent:} Detailed discussion of the procedure, potential risks, benefits, alternatives, and expected postoperative course, ensuring the patient provides informed consent.
  
  \item \textbf{Stoma Site Marking:} If there is a possibility of a temporary or permanent stoma, an enterostomal therapist will mark an optimal site on the abdomen preoperatively, considering patient comfort, clothing lines, and visibility.
\end{itemize}

\section*{Contraindications \& Precautions}
\textbf{Absolute Contraindications:}
\begin{itemize}
  \item Unresectable metastatic disease where the primary tumor is asymptomatic and surgery would not offer palliative benefit or improve survival.
  
  \item Severe, uncorrectable comorbidities that render the patient unfit for general anesthesia and major abdominal surgery (e.g., severe cardiac failure, uncontrolled sepsis, acute myocardial infarction).
  
  \item Patient refusal after thorough discussion of risks and benefits, and exploration of alternative treatment options.
\end{itemize}

\textbf{Relative Contraindications \& Precautions:}
\begin{itemize}
  \item Extensive local invasion into critical, unresectable structures (e.g., sacrum above S2, major pelvic vessels, extensive bladder or ureteral involvement) where an R0 resection is deemed impossible without unacceptable morbidity.
  
  \item Poor performance status (e.g., ECOG >3) or severe malnutrition, which significantly increases perioperative risks and impairs wound healing.
  
  \item Active, uncontrolled infection (e.g., diverticulitis, severe peritonitis) or severe inflammatory conditions that could complicate surgery or healing.
  
  \item Significant obesity (BMI >35-40 kg/m\textasciicircum{}2), which can increase technical difficulty, prolong operative time, and elevate the risk of complications, particularly with minimally invasive approaches.
  
  \item Prior extensive pelvic radiation or surgery, which can lead to dense adhesions, fibrosis, and increased risk of injury to surrounding structures, making dissection more challenging.
  
  \item Advanced age alone is not a contraindication, but careful assessment of physiological reserve, frailty, and cognitive function is paramount.
  
  \item Coagulopathy or severe bleeding disorders that cannot be adequately corrected preoperatively.
\end{itemize}

\section*{Complications \& Risks}
Rectal cancer surgery carries a range of potential risks, both general to any major surgery and specific to the procedure itself, which can significantly impact patient recovery and long-term outcomes.

\begin{itemize}
  \item \textbf{General Surgical Risks:}
    \begin{itemize}
      \item \textbf{Anesthesia Complications:} Adverse reactions to anesthetic agents, respiratory depression, cardiac events (myocardial infarction, arrhythmia), stroke, pneumonia, aspiration.
      
      \item \textbf{Bleeding:} Intraoperative hemorrhage requiring transfusion, or postoperative hematoma formation, potentially necessitating reoperation.
      
      \item \textbf{Infection:} Surgical site infection (superficial or deep), intra-abdominal abscess, sepsis, urinary tract infection, pneumonia.
      
      \item \textbf{Thromboembolic Events:} Deep vein thrombosis (DVT) and pulmonary embolism (PE), despite prophylactic measures.
      
      \item \textbf{Organ Injury:} Accidental injury to adjacent organs such as the small bowel, ureters, bladder, spleen, or major blood vessels during dissection.
    \end{itemize}
  
  \item \textbf{Procedure-Specific Risks:}
    \begin{itemize}
      \item \textbf{Anastomotic Leak:} The most serious complication, occurring when the surgical connection between bowel segments breaks down. This can lead to peritonitis, sepsis, pelvic abscess, reoperation, prolonged hospital stay, and potentially a permanent stoma. Its incidence is higher with very low anastomoses.
      
      \item \textbf{Bowel Obstruction:} Due to adhesions forming after surgery, internal hernia, or anastomotic stricture.
      
      \item \textbf{Urological Dysfunction:} Injury to pelvic autonomic nerves (hypogastric plexus) can cause neurogenic bladder, urinary retention, or chronic urinary dysfunction, requiring catheterization.
      
      \item \textbf{Sexual Dysfunction:} Nerve injury can lead to erectile dysfunction, ejaculatory dysfunction in males, and dyspareunia, vaginal dryness, or anorgasmia in females.
      
      \item \textbf{Perineal Wound Complications:} For APR, the perineal wound can be slow to heal, leading to infection, dehiscence, chronic sinus formation, or persistent pain.
      
      \item \textbf{Stoma Complications:} If a stoma is created, potential issues include parastomal hernia, prolapse, retraction, skin irritation, stenosis, high output, or necrosis.
      
      \item \textbf{Low Anterior Resection Syndrome (LARS):} A common functional complication after sphincter-preserving surgery, characterized by increased stool frequency, urgency, incontinence, clustering, and incomplete evacuation, significantly impacting quality of life.
      
      \item \textbf{Nerve Injury:} Beyond autonomic nerves, injury to somatic nerves like the obturator nerve can occur, leading to motor or sensory deficits.
      
      \item \textbf{Recurrence of Cancer:} Despite successful surgery, there is always a risk of local or distant recurrence, necessitating further treatment.
    \end{itemize}
\end{itemize}

\section*{Postoperative Recovery Timeline}
Immediately after rectal cancer surgery, the patient is transferred to the Post-Anesthesia Care Unit (PACU) for close monitoring of vital signs, pain levels, and recovery from anesthesia. Pain management is a priority, often involving patient-controlled analgesia (PCA) or epidural catheters to facilitate early mobilization. Intravenous fluids are administered, and a nasogastric tube or Foley catheter may be in place temporarily. Early ambulation is encouraged within 24 hours to prevent complications like DVT, PE, and pneumonia, and to promote bowel function return.

Over the next few days, the patient's recovery progresses. The Foley catheter is typically removed within 24-48 hours, and diet is gradually advanced from clear liquids to solids as bowel function returns (passage of flatus or stool). Drains, if placed, are monitored for output and removed when appropriate, usually when output is minimal. Patients with a new stoma receive intensive education from an enterostomal therapist on stoma care and management. Hospital stays typically range from 5 to 10 days, depending on the complexity of the surgery, the occurrence of complications, and the patient's overall recovery trajectory. Upon discharge, patients receive detailed instructions on wound care, pain medication, activity restrictions (e.g., no heavy lifting for 6-8 weeks), and warning signs to report. Long-term recovery involves potential stoma reversal (if temporary), adjuvant chemotherapy or radiation based on pathology results, and ongoing surveillance for recurrence. Full recovery of energy levels and return to normal activities can take several weeks to months.

\section*{Surgical Findings \& Pathology}
During rectal cancer surgery, the surgeon meticulously documents intraoperative findings, including the exact location and size of the tumor, its relationship to surrounding structures, the presence of any suspicious lymph nodes, and the integrity of the mesorectal fascia during dissection. Any unexpected findings, such as metastatic disease, are also noted. The resected specimen, comprising the rectum and its entire mesorectum, is then sent for detailed pathological examination.

The pathology report is the most critical "result" of rectal cancer surgery, providing essential information that guides further treatment decisions and predicts prognosis. Key elements include:

\begin{itemize}
  \item \textbf{Tumor Type and Grade:} Confirms adenocarcinoma and assesses its differentiation (well, moderately, poorly differentiated), which correlates with aggressiveness.
  
  \item \textbf{Tumor Stage (pTNM):}
    \begin{itemize}
      \item \textbf{pT (Pathological Tumor):} Depth of invasion into the rectal wall and surrounding tissues (e.g., T1: submucosa, T2: muscularis propria, T3: through muscularis propria into perirectal fat, T4: invades adjacent organs/peritoneum).
      
      \item \textbf{pN (Pathological Nodes):} Number of regional lymph nodes involved with cancer. An adequate lymph node harvest (typically >12 nodes) is crucial for accurate staging.
      
      \item \textbf{pM (Pathological Metastasis):} Confirmation of no distant metastases in the resected specimen (though distant staging is primarily done preoperatively).
    \end{itemize}
  
  \item \textbf{Resection Margins:} Crucially, whether the surgical margins are clear of cancer (R0 resection), microscopically positive (R1), or macroscopically positive (R2). An R0 resection is the primary goal for curative intent. The circumferential resection margin (CRM), the shortest distance from the tumor to the radially cut edge of the mesorectum, is particularly important for rectal cancer, with a margin of $\geq$1 mm considered clear.
  
  \item \textbf{Extramural Vascular Invasion (EMVI) and Perineural Invasion (PNI):} Presence of cancer cells in blood vessels or nerves outside the rectal wall, which are adverse prognostic factors indicating a higher risk of recurrence.
  
  \item \textbf{Response to Neoadjuvant Therapy:} For patients who received preoperative treatment, the pathology report will assess the degree of tumor regression (e.g., Mandard tumor regression grade), which is an important prognostic indicator.
\end{itemize}

These findings collectively determine the patient's pathological stage, risk of recurrence, and the necessity for adjuvant (postoperative) chemotherapy or radiation therapy.

\section*{Expected Postoperative Course}
A normal, successful postoperative course following rectal cancer surgery signifies both effective oncological clearance and a smooth functional recovery. Patients can expect initial pain managed by analgesia, which gradually subsides over days to weeks. Bowel function typically returns within 3-5 days, marked by the passage of flatus and then stool, allowing for a progressive diet advancement. Surgical wounds should heal without signs of infection or dehiscence. For those with a temporary stoma, learning stoma care and management becomes part of the normal course. While some degree of fatigue is common for several weeks, patients should experience a gradual return of energy and ability to perform daily activities. For sphincter-preserving procedures, a normal course implies the patient is able to manage their bowel function with an acceptable quality of life, even if some degree of Low Anterior Resection Syndrome (LARS) is present and manageable with dietary adjustments or medications. The ultimate goal is a return to a functional baseline, allowing for subsequent adjuvant therapies if indicated, and long-term surveillance.

\section*{Postoperative Care \& Follow-up}
Postoperative care and long-term follow-up are critical for monitoring recovery, detecting recurrence, and managing potential long-term complications.

\begin{itemize}
  \item \textbf{Postoperative Clinic Visits:} Regular visits with the surgeon and oncologist are scheduled to monitor wound healing, review the final pathology report, discuss adjuvant therapy plans, and address any symptoms or concerns.
  
  \item \textbf{Adjuvant Therapy:} If indicated by the pathology report (e.g., for Stage III disease or high-risk Stage II), adjuvant chemotherapy and/or radiation therapy will be initiated, typically within 4-8 weeks post-surgery.
  
  \item \textbf{Surveillance Imaging:} Regular CT scans of the chest, abdomen, and pelvis are typically performed every 6-12 months for 3-5 years to detect distant metastases. Pelvic MRI may be used for local recurrence surveillance, especially for patients with high-risk features.
  
  \item \textbf{Colonoscopy:} A surveillance colonoscopy is usually performed one year after surgery (or one year after stoma reversal) to check for anastomotic recurrence, new polyps, or metachronous cancers, then typically every 3-5 years, depending on findings and risk factors.
  
  \item \textbf{CEA Tumor Marker:} Blood tests for carcinoembryonic antigen (CEA) levels are monitored every 3-6 months for 3-5 years, as rising levels can indicate recurrence.
  
  \item \textbf{Rehabilitation and Support:} Patients may benefit from pelvic floor physical therapy for LARS symptoms, and support groups for stoma management or cancer survivorship. Nutritional counseling is also often provided.
  
  \item \textbf{Activity Resumption:} Gradual return to normal activities, with specific restrictions on heavy lifting for several weeks to months to prevent incisional hernias.
\end{itemize}

Patients are educated on warning signs such as severe abdominal pain, persistent fever, uncontrolled nausea/vomiting, significant changes in bowel habits, unexplained weight loss, or new rectal bleeding, which warrant immediate medical attention.

\section*{Epidemiology}
Rectal cancer accounts for approximately one-third of all colorectal cancer diagnoses, making it a significant public health concern globally. Incidence rates vary geographically, with higher rates observed in developed countries, often linked to Westernized diets and lifestyles. While traditionally a disease of older adults, typically affecting individuals over 50 years of age, there has been an alarming increase in the incidence of early-onset rectal cancer (diagnosed before age 50) in recent decades, for reasons that are not yet fully understood. The disease shows a slight male predominance. The most common indication for rectal cancer surgery is adenocarcinoma, which constitutes over 95\% of all rectal malignancies. Despite advancements in surgical techniques and neoadjuvant therapies, rectal cancer surgery is associated with significant morbidity, with complication rates varying based on the procedure type, patient factors, and surgical expertise. Mortality rates have steadily declined over the past few decades due to improved surgical expertise, perioperative care, and multidisciplinary management, but still represent a serious risk, particularly in elderly or frail patients.

\section*{Outcomes \& Prognosis}
The outcomes and prognosis following rectal cancer surgery are highly dependent on the stage of the disease at diagnosis, the quality of the surgical resection (particularly the circumferential resection margin), and the patient's response to neoadjuvant and adjuvant therapies. With modern Total Mesorectal Excision (TME) techniques, local recurrence rates have dramatically decreased from historical highs of 30-50\% to 5-10\%, as demonstrated by landmark studies like the Dutch TME trial. Five-year survival rates vary significantly by pathological stage: Stage I tumors typically have survival rates exceeding 90\%, while Stage II tumors range from 70-80\%, and Stage III tumors from 50-70\%.

Functional outcomes are a major consideration, particularly after sphincter-preserving procedures. Low Anterior Resection Syndrome (LARS), characterized by altered bowel function (increased stool frequency, urgency, incontinence, clustering, and incomplete evacuation), affects a significant proportion of patients (up to 80\% in some series), impacting their quality of life. Key prognostic factors include the pathological TNM stage, the status of the circumferential resection margin (CRM), the number of positive lymph nodes, the presence of extramural vascular invasion (EMVI) or perineural invasion (PNI), and the degree of tumor regression after neoadjuvant therapy. Patients with an R0 resection and a good pathological response to neoadjuvant therapy generally have the best prognosis. While local control is excellent with TME, distant recurrence remains a primary concern, often occurring in the liver or lungs. Ongoing surveillance is crucial for early detection and management of recurrence, which can significantly influence long-term survival.

\section*{References}
\begin{enumerate}
  \item MedlinePlus. Rectal cancer surgery. U.S. National Library of Medicine; 2024.
  
  \item Schwartz's Principles of Surgery. 12th ed. McGraw-Hill Education; 2023.
  
  \item NCCN Clinical Practice Guidelines in Oncology: Rectal Cancer, Version 1.2024. National Comprehensive Cancer Network; 2024.
  
  \item Sabiston Textbook of Surgery. 22nd ed. Elsevier; 2022.
\end{enumerate}

\clearpage